\documentclass[]{../../../tex/classes/styledArticle}
\usepackage{../../../tex/packages/hebrewSupport}


\renewcommand\lg{\lambda}
\newcommand\ag{\alpha}
\newcommand\cm{\mathrm{cm}}
\newcommand\lnsmm{\frac{\mathrm{lines}}{\mathrm{mm}}}
\newcommand\mrm[1]{\mathrm{#1}}
\newcommand\sof[1]{|#1|}
\newcommand\cl[1]{(#1)}
\newcommand\N{\mathbb{N}}
\author{פרץ קורן שפיר בוילר}
\title{מעבדה ראשונה בפיזיקה}
\begin{document}
	\maketitle
	
	\section{שאלות פתיחה}
	\subsection{שאלה 1}
	נוכיח ש־$\Delta x = \frac{L \lg}{d}$
	\begin{proof}
		תחת ההנחה ש־$L \gg d$, ניתן להשתמש בטריגו על שני משולשים כדי להבחין בשתי הנוסחאות: 
		\begin{align*}
			(1)&\quad\quad\sin \theta_k = \frac{k \lg}{d} \\
			(2)&\quad\quad\tan \theta_k = \frac{x_k}{L}
		\end{align*}
		באמצעות קירוב זווית קטנה, נסיק ש־$\sin \theta_k \approx \tan \theta_k$. לכן: 
		\[ \frac{k\lg}{d} = \sin \theta_k \approx \tan \theta_k = \frac{x_k}{L} \]
		נתבונן בהפרש $\Delta x_k$ בין $x_k$ ו־$x_{k +1}$: 
		\[ \Delta x_k = x_{k + 1} - x_k = (k + 1)\frac{L\lg}{d} - k \cdot \frac{L\lg}{d} = \frac{L\lg }{d}\cdot\underbrace{((k + 1) - k)}_1 = \frac{L\lg}{d}  \]
		נבחין שבעבור זוויות קטנות, אין תלות ב־$k$. סה''כ $\Delta x = \frac{L \lg }{d}$ כנדרש. 
	\end{proof}
	
	\subsection{שאלה 2}
	נסביר מהם שני ההבדלים בין תמונת ההתאבדות בין 2 סדקים צרים לבין תבנית עקיפה כאשר האור עובר דרך סדק יחיד. 
	
	\textbf{תשובה: }בחריץ יחיד תמונת ההתאבכות מורכבת מפס אמצעי בעל $\Delta x$ שגדול פי 2 מזה של הקווים בסדרים הבאים. זאת לעומת הקווים האחידים בעוביים עבור שני חריצים. בנוסף, בחריץ יחיד העוצמה משתנה ככל שמתרחקים בסדרים, בעוד שבשני חריצים העוצמה נשארת קבועה בסדרים הראשונים. 
%	בחריץ יחיד תמונת ההתאבכות מורכבת מפס אמצעי בעל $\Delta x$ שווה לפי 2 מזה של הקווים עבור הסדרי םהבאים ולעומת הקווים האחידים בשני חריצים. 
	\subsection{שאלה 3}
	נגדיר את המושגים הבאים. 
	\begin{itemize}
		\item \textbf{קבוע שריג: }$N^*$, כמות החריצים ליחידת אורך (לרוב, כמות החריצים ל־$\mathrm{mm}$, כלומר נמדד ביחידות של $\frac{1}{\mathrm{mm}}$), ישנו קשר הפוך $N^* = \frac{1}{d}$ לפי הגדרה. 
		\item \textbf{כושר הפרדה: }
		\item \textbf{מקורות אור קוהרנטיים: }מקורות אור בעלי זמן מחזור ואורך גל זהים
	\end{itemize}
	
	\subsection{שאלה 4}
	נוכיח שכאשר מועבר אור לבן דרך סריג עקיפה, רק הסדר הראשון משני צידי האנך המרכזי מתקבל באופן מלא וברור. 
	\begin{proof}
		???
	\end{proof}
	
	
	\section{הניסוי}
	לאורך כל ניסוי $L = 10\cm$. 
	\subsection{שאלה 1}
	נוכיח ש־$\Delta x = L \lg N^*$
	\begin{proof}
		ידוע מנוסחה (1): 
		\[ \sin a_k = k \lg N^* \]
		ידוע מנוסחה (2): 
		\[ \tan a_k = \frac{x_k}{L} \]
		באמצעות קירוב זווית קטנה, נסיק ש־$\sin a_k \approx \tan a_k$. לכן: 
		\[ k\lg N^* = \sin a_k \approx \tan a_k = \frac{x_k}{L} \]
		דהיינו $\Delta x_k = k L \lg N^*$. 
		נתבונן בהפרש $\Delta x_k$ בין $x_k$ ו־$x_{k +1}$: 
		\[ \Delta x_k = x_{k + 1} - x_k = (k + 1)L\lg N^* - kL\lg N^* = L\lg N^*\cdot\underbrace{((k + 1) - k)}_1 = L\lg N^*  \]
		נבחין שבעבור זוויות קטנות, אין תלות ב־$k$. סה''כ $\Delta x = L \lg N^*$ כנדרש. 
	\end{proof}
	
	\subsection{שאלה 2}
	בעבור סריג של $1000\frac{\mathrm{lines}}{\mathrm{mm}}$: 
	\begin{enumerate}[(A)]
		\item המרחק $\Delta x$ מהנקודה המרכזית לנקודה שמימין לה: $\bm{8.5\cm}$
		\item המרחק $\Delta x$ מהנקודה המרכזית לנקודה שמשאל לה: $\bm{8.4\cm}$
		\item מדוע עדיף לבצע מדידות משני צדי הנקודה המרכזית? 
		
		\textbf{תשובה: }עדיף לבצע מדידות משני צדי הנקודה המרכזית כי 
		
		\item ממוצע שתי המדידות: $\bm{8.4\cm}$
	\end{enumerate}
	
	\subsection{שאלה 3}
	בעבור סריג של $N^*_1 := 500\frac{\mathrm{lines}}{\mathrm{mm}}$: 
	\begin{enumerate}[(A)]
		\item המרחק בין שתי נקודות אור מסדר ראשון: $\bm{6.9\cm}$, ולכן $\Delta x \approx \bm{3.45\cm}$. 
		\item מדוע שיטה זו מדויקת יותר: 
		
		\textbf{תשובה: }מדידה של מרחק גדול יותר מאפשרת סטייה יחסית קטנה יותר. 
		
%		אנו חוסכים את חוסר הדיוק של הנייר המילימטרי כאשר מודדים פעמיים את 0. השגיאה היחסית של $\pm 0.5\mathrm{mm}$ לכל מדידה קטנה יותר כשמודדים מרחק גדול. 

		\item אם היינו מציבים סריג עם צפיפות $N^*_2 := 300\frac{\mathrm{lines}}{\mathrm{mm}}$, אזי: 
		
		\begin{align*}
			\begin{cases}
				3.45 = &\!\!\!\!\Delta x_1 = L \cdot \lg \cdot N^*_1 \\
				&\!\!\!\!\Delta x_2 = L \cdot \lg \cdot N_2^*
			\end{cases}
		\end{align*}
		נחלק את שתי המשוואות ונקבל: 
		\[ \frac{\Delta x_2}{3.45} = \frac{\Delta x_2}{\Delta x_1} = \frac{\cancel{L \cdot \lg} \cdot N^*_2}{\cancel{L \cdot \lg} \cdot N_1^*} = \frac{300}{1000} = 0.3 \]
		סה''כ: 
		\[ \Delta x_2 = 3.45 \cdot 0.3 = 1.035\cm \]
		\item לאחר סיבוב הסריג והסרגל ב־$90^{\circ}$, באיזו תבנית נבחין? 
		
		\textbf{תשובה: }התבנית הסתובבה אנכית ב־$90^{\circ}$ מעלות, כי החריצים בסריג הסתובבו ב־$90^{\circ}$. 
	\end{enumerate}
	
	\subsection{שאלה 4}
	נמקם את הסריג עם $N^* = 190\frac{\mathrm{lines}}{\mathrm{mm}}$. 
	\begin{enumerate}[(A)]
		\item נתאר את תבנית ההתאבכות. 
		
		\textbf{תשובה: }נבחין בתבנית התאבכות של ''מטריצה`` של נקודות, שבסביבת 0 נראים כמו מעין ריבועים. 
		\item מדוע מתקבלת התבנית הזו: 
		
		\textbf{תשובה: }לדעתנו התבנית מתקבלת כי ישנם הן קווים אופקיים והן קווים אנכיים בסריג. 
		
		\item נמצא את $\Delta x$ בין שתי נקודות אופקיות: $\bm{2.6\cm}$
		\item נמצא את $\Delta x$ בין שתי נקודות אנכיות: $\bm{2.5\cm}$
		\item האם קבוע־הסריג בממד האופקי זהה לזה שבממד האנכי? 
		
		לדעתנו, קבוע הסריג בשני הממדים זהה. זאת כי קבוע הסריג פורפורציונלי ביחס ישר ל־$\Delta x$, והם יצאו זהים בסעיפים (ב) ו־(ג) עד כדי שגיאת מדידה. 
	\end{enumerate}
	
	\subsection{שאלה 5}
	נבצע את הניסוי עם תקליטור CD. צפיפות החריצים שלו $N^* = 680\lnsmm$. 
	
	נבחין ש־$\Delta x= 4.85\cm$, כמו בסעיף א' בשאלה 3. 
	
	\subsection{שאלה 6}
	ניעזר בנתונים מהשאלות השונות כדי לחרבש טבלה. 
	
	TODO
	
	\begin{center}
		\begin{tabular}{c|c|c|c|c|c} 
			5 & 4ב' & 3ג & 3א & 2ד & המשתנה \\ \hline
			680  & 190 & 300   & 500  & 1000 & $N^* \left[{\frac{1}{\mathrm{mm}}}\right]$ \\ \hline
			48.5 & 26  & 10.35 & 34.5 & 84   & $\Delta x \left[{{\mathrm{mm}}}\right]$ \\ \hline
			&  &  &  & & ${\tan \ag_1}$ \\ \hline
			&  &  &  & & ${\ag_1}$ \\ \hline
			&  &  &  & & ${\sin \ag_1}$ \\ 
		\end{tabular}
	\end{center}
	
	\subsection{שאלה 7}
	TODO
	
	\begin{enumerate}[(A)]
		\item נסביר מדוע הקשר בין $\sin \ag_1$ ו־$N^*$ הוא לינארי. 
		
		\textbf{תשובה: }נבחין במשוואה (1) שברקע העיוני, $\sin \ag_1 = N^* \lg k$. משום ש־$\lg$ ו־$k$ לא תלויים ב־$\ag_1$, אזי מהגדרה $\sin\ag_1 \propto N^*$ כנדרש. 
		
		\item נקבע מהו המשתנה התלוי ומהו המשתנה הבלתי־תלוי. 
		
		\item נסרטט על הנייר המילימטרי (ידוע בכינויו ''אקסל``) דיאגמת פיזור של המשתנה התלוי כפונקציה של המשתנה הבלתי־תלוי, ע''פ התוצעות בשאלה 6. 
		\item נעביר קו מגמה בתרשים. 
		\item נחשב (כלומר אקסל יחשב) את שיפוע קו המגמה. יחידות המידה שלו הן
		\item נמצא את $\lg$. 
		
		\textbf{תשובה: }
	\end{enumerate}
	
	\section{שאלות נוספות}
	\subsection{שאלה 2}
	תלמיד מבצע ניסוי בו הוא משתמש בסריג עקיפה עם $10^{4}$ קווים לס''מ, ובמסך ארוך מאוד המוצב במקביל לסריג במרחק $L = 2.5\mathrm{m}$ ממנו. בהתחלת בניסוי התלמיד השתמש בפנס לייזר, הפולט אור מונוכרומטי ב־$\lg = 400\mathrm{nm}$. אור הלייזר פוגע בסריג במאונך לו. 
	\begin{enumerate}[(A)]
		\item נחשב את מרחק המקסימום מסדר ראשון מהאנך האמצעי. 
		
		\textbf{תשובה: }נמיר יחידות מידה ונגלה ש־$N^* = 10^{6}\frac{1}{\mrm{m}}, \ \lg = 4 \cdot 10^{-7}\mrm{m}, L = 2.5\mrm{m}$. נציב בנוסחה: 
		\[ \Delta x_1 = L \cdot \frac{\lg}{d} = {2.5 \cdot 4 \cdot 10^{-7} \cdot 10^{6}} = \bm{1 \mrm{m}} \]
		
		\item נסמן את מספר קווי המקסימום ב־$N$. ידוע $n_{\max} \le \frac{d}{\lg}$, כלומר: 
		\begin{align*}
			N &= 2n_{\max} + 1 \\
			&= 2 \cdot \left\lfloor \frac{d}{\lg} \right\rfloor + 1 \\
			&= 2 \cdot \left\lfloor \frac{1}{N^*\lg} \right\rfloor + 1 \\
			&= 2 \cdot \left\lfloor \frac{1}{10^{6} \cdot 4 \cdot 10^{-7}} \right\rfloor + 1 = 2 \cdot 2 + 1 = \bm{5}
		\end{align*}
		\item אם אותה המערכת תהיה בתוך שמן בעל מקדם שבירה $1.6$, נבחין שמספר קווי המקסימום יהיה 
		
		\textbf{תשובה. }לעשות
		
		\item התלמיד מנסה לקבל במרחק $2.5\mrm{m}$ תבנית התאבכות דרך סריג עם קרני X (אורך גל $0.1\mrm{nm} \le \lg \le 1\mrm{mm}$), אך אינו מצליח. נסביר את הסיבה. 
		
		\textbf{תשובה: }בהתאם לנוסחה הקודמת, נקבל שכמות קווי המקסימום תהיה: 
		\[ N = 2 \cdot \left\lfloor \frac{1}{N^*\lg} \right\rfloor + 1 \ge 2 \left\lfloor \frac{1}{10^{6} \cdot 1 \cdot 10^{-9}} \right\rfloor = 2001 \]
		כלומר, אותו האור מתחלק על 2001 קווי מקסימום על המסך. זוהי כמות גדולה של קווי מקסימום, מעבר לכושר הראייה של התלמיד, ולכן הוא לא יצליח לראות תבנית התאבכות. 
	\end{enumerate}
	\subsection{שאלה 3}
	תלמיד התבקש להרכיב מכשיר להפרדת אורכי גל שונים של מקור גלים אלקטרומגנטיים. נתונים אורכי גל $\lg_1 = 5000\mrm{A}^{\circ} = 5 \cdot 10^{-7} \mrm{m}$ ו־$\lg_2 = 6600\mrm{A}^{\circ} = 6.6\cdot 10^{-7} \mrm{m}$. מרחק שני סדקים בשקופית הפלסטיק שלו הוא $d = 0.006\mrm{mm} = 6\cdot 10^{-6} \mrm{m}$. במרחק $L = 5\mrm{m}$ מהשקופית ישנו מסך. 
	
	\begin{enumerate}[(A)]
		\item נתבונן בסרטוט הנתון בשאלון. נבחין שנקודה $A$ על האנך המרכזי, ונקודה $B$ קרובה אליה מנקודה $C$. שתיהן מופיעות בקו המקסימום מסדר ראשון של הגלים המתאימים להם. 
		
		נסמן ב־$\ag_i^{(1)}$ את קו המקסימום מסדר $i$ של $\lg_1$ וב־$\lg_i^{(2)}$ את קו המקסימום מסדר $i$ של $\lg_2$. אז: 
		\begin{alignat*}{9}
			\sin \ag_1^{(1)} &= \frac{i\lg_1}{d} = \frac{1 \cdot 5 \cdot 10^{-7}}{6 \cdot 10^{-6}} &&= 8.33\cdot 10^{-2} && \quad\implies \quad \ag_1^{(1)} = \arcsin(8.33 \cdot 10^{-2}) = 4.78^{\circ} \\
			\sin \ag_1^{(2)} &= \frac{i\lg_2}{d} = \frac{1 \cdot 6.6 \cdot 10^{-7}}{6 \cdot 10^{-6}} &&=1.1\cdot 10^{-1} && \quad\implies \quad \ag_1^{(2)} = \arcsin(1.1 \cdot 10^{-1}) = 6.32^{\circ} \\
		\end{alignat*}
		נבחין ש־$\ag_1^{(2)} > \ag_1^{(1)}$ ולכן הנקודה הרחוקה יותר (היא נקודה C) שייכת ל־$\lg_2$ בעוד הנקודה הקרובה יותר (היא נקודה B) שייכת ל־$\lg_1$. 
		\item נמצא את המרחק בין הנקודות B ו־C. 
		
		\textbf{תשובה: }מנוסחת יאנג וקירובי זווית קטנה, ידוע ש־: 
		\begin{alignat*}{9}
			\Delta x^{(1)} &= \frac{\lg_1 \cdot L}{d} = \frac{\sin\ag_1^{(1)}}{L} &&=  \frac{8.33\cdot 10^{-2}}{5} = 1.67 \cdot 10^{-2}\mrm{m} \\
			\Delta x^{(2)} &= \frac{\lg_2 \cdot L}{d} = \frac{\sin\ag_1^{(2)}}{L} &&=  \frac{1.1\cdot 10^{-1}}{5} = 2.2 \cdot 10^{-2}\mrm{m}
		\end{alignat*}
		לכן המרחק הוא: 
		\[ \sof{\Delta x^{(2)} - \Delta x^{(1)}} = \cl{2.2 - 1.67} \cdot 10^{-2} = 0.53 \cdot 10^{-2}\mrm{m} \]
		
		\item נחפש את הסדר הנמוך ביותר פרט לסדר $0$ עבורו הקווים מתלכדים (אם קיים כזה). לשם כך, נחפש $i, j \in \N_+$ כך ש־$\ag^{(1)}_i = \ag^{(2)}_j$. באופן שקול נדרוש ש־$\sin \ag^{(1)}_i = \sin \ag^{(2)}_j$, כלומר: 
		\begin{align*}
			\sin \ag^{(1)}_i &= \sin \ag^{(2)}_j \\
			\frac{i \lg_1}{d} &= \frac{j\lg_2}{d} \\
			i \cdot \frac{5 \cdot 10^{-7}}{6 \cdot 10^{-6}} &= j \cdot j \cdot \frac{6.6 \cdot 10^{-7}}{6 \cdot 10^{-6}} \\ 
			i \cdot 8.33\cdot 10^{-2} &= j 1.1\cdot 10^{-1} \\
			i \cdot 833 &= j 1100
		\end{align*}
		נבחין שמתקיים שוויון (עד לכדי שגיאת חישוב זניחה) בעבור $i = 4, j = 3$. הפתרון תקין כי $j \cdot 1.1 \cdot 10^{-1} \in (0, 1)$. סה''כ הקו מסדר $3$ של $\lg_1$ המינימלי מסדר לא טרוויאלי שמתלכד עם קו מקסימום של $\lg_2$. 
		
		\item נניח שהמסך גדול לאין־שיעור וכן שתנאי הריאות אופטימליים. נמצא את כמות קווי המקסימום של $\lg_1$. 
		
		\textbf{תשובה. }
	\end{enumerate}
	
	
	\subsection{שאלה 4}
	TODO
	\subsection{שאלה 5}
	TODO
	
	
	
	
%	\ndoc
\end{document}